\paragraph{\label{sec:introduction}Introduction --}

Astrophysical and cosmological observations have provided compelling evidence for the existence of dark matter (DM) that would constitute most of the matter content of the Universe~\cite{1937ApJ,1970ApJ...159..379R,Clowe_2006,Planck2018VI}.
Weakly Interacting Massive Particles (WIMPs) are among the best-motivated candidates for dark matter, arising naturally in many theories beyond the Standard Model.
Numerous direct-detection experiments have searched for interactions between dark matter particles and ordinary matter, with experiments such as LUX-ZEPLIN (LZ) placing stringent constraints on many leading theoretical models~\cite{LZ:SR3_WS2024, XenonnT:WIMP-SI-2025,PandaX4T:WIMP-2025,DEAP3600:2026_WIMP,DarkSide-50:2018_532days,SuperCDMS:2018_WIMP,EDELWEISS:2019_WIMP,CRESST:2019_DM,PICO:2019_DM}.

WIMP searches for particles with mass  $\gtrsim10$~GeV/c$^2$ typically focus on spin-independent (SI) and spin-dependent~(SD) interactions generating $\lesssim100$~keV nuclear recoils~\cite{Lewin:1995rx, Goodman:1984dc, Griest:1988ma}.
Any momentum dependence in these interactions is assumed to be suppressed as the inverse of the momentum transfer is much larger than the size of the nucleon. 
Allowing for momentum-dependent WIMP-nucleon couplings can result in predicted recoil spectra that differ from the SI and SD cases, including non-negligible contributions at higher recoil energies~\cite{Chang:2009yt,Feldstein:2009tr}. 
In this Letter, we report the latest results from the LZ experiment in a nuclear recoil energy range up to 270~keV.


